% Part: computability
% Chapter: recursive-functions
% Section: normal-form

\documentclass[../../../include/open-logic-section]{subfiles}

\begin{document}

\olfileid[pt]{cmp}{rec}{nft}
\olsection{O teorema da forma normal}

\begin{thm}[Teorema da forma normal de Kleene]
\ollabel{thm:kleene-nf}
Existe uma relação recursiva primitiva $T(e, x, s)$ e uma função
recursiva primitiva $U(s)$ com a seguinte propriedade: se $f$ é uma
função recursiva parcial qualquer, então, para algum~$e$,
\[
f(x) \simeq U(\umin{s}{T(e, x, s)})
\]
para todo $x$.
\end{thm}

\begin{explain}
A demonstração do teorema da forma normal é complexa, mas a ideia
básica é simples. Toda função recursiva parcial tem um \emph{índice}~$e$,
intuitivamente, um número que codifica seu programa ou sua definição.
Se $f(x) \fdefined$, a computação pode ser registrada sistematicamente
e codificada por algum número~$s$, e o fato de que $s$ codifica a
computação de~$f$ na entrada~$x$ pode ser verificado recursivamente de
modo primitivo usando apenas $x$ e a definição~$e$. Consequentemente,
a relação~$T$, ``a função de índice~$e$ tem uma computação para a
entrada~$x$, e $s$ codifica essa computação'', é recursiva primitiva.
Dado o registro completo da computação~$s$, o ``resultado'' de~$s$ é
o valor de~$f(x)$, que também pode ser obtido de~$s$ recursivamente de
modo primitivo.

O teorema da forma normal mostra que basta uma única busca ilimitada
para definir qualquer função recursiva parcial. Basicamente, podemos
percorrer todos os números até encontrar um que codifique uma computação
da função de índice~$e$ para a entrada~$x$. Podemos usar os números~$e$
como ``nomes'' de funções recursivas parciais e escrever $\cfind{e}$
para a função~$f$ definida pela equação do teorema. Observe que uma
função recursiva parcial pode ter mais de um índice---na verdade, toda
função recursiva parcial tem infinitos índices.
\end{explain}
\end{document}

